% Осень26!B5; нумерация внутри столбца без пустых строк.
% Источник: лекция 4, §14 (строка 454)
\definitionticket{03}{Локальный максимум, минимум, экстремум}
\label{definition:03}

\subsubsection*{Определения}

Пусть \(a\in\mathrm{Int}\,D\).

\begin{itemize}
\tightlist
\item
  \textbf{Локальный максимум:} \(\exists U(a)\subset D:\ \forall x\in U(a)\ f(x)\le f(a)\) (нестрогое --- поблизости могут быть точки с тем же значением).
\item
  \textbf{Строгий локальный максимум:} \(\forall x\in\dot U(a)\) (проколотая окрестность) \(f(x)<f(a)\). Вставка прокола в предыдущее определение смысла не меняет, а здесь --- даёт именно строгий знак.
\item
  \textbf{Локальный минимум} / \textbf{строгий локальный минимум}: аналогично с \(\ge\) (\(>\)).
\item
  \textbf{Локальный экстремум} --- максимум или минимум (строгий --- аналогично).
\end{itemize}

(Шутка лектора про «оценку получил, но не понял, положительная или отрицательная» --- про ситуацию «то ли максимум, то ли минимум, но мы ещё не разобрались».)
